\chapter{Conclusion}

This thesis set out to address a long-standing challenge in Quantitative Finance: capturing the stylized facts of financial time series: heavy tails, volatility clustering, and the leverage effect, within a generative framework that is both expressive and tractable. Two classical baselines, GARCH-X and Merton-X, were first calibrated as a point of reference (Ch.\ \ref{ch:experiments_results}): well understood, interpretable, and partially effective at reproducing the bulk of the distribution, but structurally unable to capture the leverage effect and other distributional properties. This established the gap that a deep generative approach would need to close.

Phase~1 replicated the CSDI-based diffusion architecture of Kim et al.\ \cite{kim2025diffusion}, whose attention-based backbone descends from a lineage of gated residual architectures, PixelCNN, WaveNet, and DiffWave, across the VE, VP, and GBM SDE families and several noise schedules. This phase confirmed that the architecture can learn the bulk of the return distribution and recover, at reduced scale, key stylized facts such as the leverage effect, while also clear limitations arose in the tails and higher moments, most severely for the GBM family. These findings both validated the chosen architecture and motivated the extension pursued in Phase~2.

The central contribution of this thesis, EDM-CSDI, addressed these limitations directly by embedding the same CSDI backbone within the Elucidated Diffusion Model framework of Karras et al.\ \cite{karras2022elucidating} and introducing conditioning on same-day Open, High, and Low log-returns. As shown in Ch.\ \ref{ch:experiments_results} and discussed in Ch.\ \ref{ch:discussion}, this combination substantially narrowed the gap left open by Phase~1: heavy tails, volatility clustering, and the leverage effect were reproduced far more faithfully, and the conditioning information translated into sharper, better-informed probabilistic forecasts than an unconditional model could provide.

Taken together, these results answer both research questions posed at the beginning: a CSDI-based diffusion architecture is indeed able to learn the stylized facts of financial time series, and embedding it within the EDM framework while conditioning on readily available intraday information meaningfully improves its ability to do so. Beyond the empirical findings, this thesis delivers a clarified and generalized EDM-CSDI design that is mathematically transparent and replicable, offering a solid foundation for the open questions and extensions outlined in Ch.\ \ref{ch:discussion}. Overall, the work presented here is an encouraging step toward generative models of financial time series that can be relied upon for future expansion in scenario analysis, stress testing, and risk-aware decision making.
